\documentclass[11pt,a4paper]{article}

\usepackage{fontspec}
\IfFontExistsTF{Times New Roman}{\setmainfont{Times New Roman}}{\setmainfont{Tinos}}
\IfFontExistsTF{Consolas}{\setmonofont{Consolas}[Scale=0.92]}{\setmonofont{Liberation Mono}[Scale=0.92]}
\usepackage[top=1.45cm,bottom=1.55cm,left=1.7cm,right=1.7cm]{geometry}
\usepackage{xcolor}

% --- OFFICIAL EXAM HEADER COLORS ---
\definecolor{ForestGreen}{RGB}{22,100,70}
\definecolor{SoftGreen}{RGB}{229,244,235}
\definecolor{CodeBgHeader}{RGB}{248,250,249}
\definecolor{DarkGrayHeader}{RGB}{55,55,55}
\usepackage{amsmath,amssymb}
\usepackage{tcolorbox}
\tcbuselibrary{skins}
\usepackage{booktabs}
\usepackage{tabularx}
\usepackage{array}
\usepackage{listings}
\usepackage{fancyhdr}
\usepackage{lastpage}
\usepackage{enumitem}
\usepackage{titlesec}
\usepackage{hyperref}

% --- COMPACT EXAM LAYOUT ---
\setlength{\parindent}{0pt}
\setlength{\parskip}{0pt}
\setlist{topsep=1.5pt,itemsep=1pt,parsep=0pt,partopsep=0pt}
\setlist[enumerate,2]{topsep=0.5pt,itemsep=0pt,parsep=0pt,partopsep=0pt}
\titlespacing*{\section}{0pt}{4pt}{2pt}
\titlespacing*{\subsection}{0pt}{3pt}{1.5pt}


\definecolor{primaryblue}{HTML}{0F3460}
\definecolor{secondaryteal}{HTML}{0D9488}
\definecolor{accentgold}{HTML}{D97706}
\definecolor{darkgray}{HTML}{1F2937}
\definecolor{lightgray}{HTML}{F3F4F6}
\definecolor{boxbg}{HTML}{F8FAFC}
\definecolor{borderblue}{HTML}{93C5FD}
\definecolor{pyblue}{HTML}{1F77B4}
\definecolor{pygreen}{HTML}{15803D}
\definecolor{pyorange}{HTML}{C2410C}

\lstdefinestyle{pythonstyle}{
  language=Python,
  basicstyle=\ttfamily\small,
  keywordstyle=\color{pyblue}\bfseries,
  commentstyle=\color{pygreen}\itshape,
  stringstyle=\color{pyorange},
  showstringspaces=false,
  columns=fullflexible,
  keepspaces=true,
  frame=single,
  rulecolor=\color{gray!35},
  backgroundcolor=\color{boxbg},
  breaklines=true,
  tabsize=4,
  aboveskip=1.5pt,
  belowskip=1.5pt
}
\lstset{style=pythonstyle}

\hypersetup{
  colorlinks=true,
  linkcolor=primaryblue,
  urlcolor=secondaryteal,
  pdftitle={DSAI1003 - Exam Part 01: Introduction, Computers, Algorithms and Basic I/O},
  pdfauthor={Dr. Vu Duc Minh - National Economics University (NEU)}
}

\pagestyle{fancy}
\fancyhf{}
\setlength{\headheight}{14pt}
\fancyhead[L]{\footnotesize\textbf{DSAI1003 -- Fundamental Programming Concepts in Python}}
\fancyhead[R]{\footnotesize\textbf{Midterm Progress Exam: Part 01 -- 60 Minutes}}
\fancyfoot[L]{\scriptsize National Economics University -- College of Technology}
\fancyfoot[R]{\footnotesize Page \thepage\ of 6}
\renewcommand{\headrulewidth}{0.4pt}
\renewcommand{\footrulewidth}{0.3pt}

\titleformat{\section}{\color{primaryblue}\Large\bfseries}{\thesection}{1em}{}[\titlerule]
\titleformat{\subsection}{\color{secondaryteal}\normalsize\bfseries}{\thesubsection}{1em}{}

\newtcolorbox{answerbox}[2][]{
  enhanced,
  colback=white,
  colframe=borderblue,
  arc=1.2mm,
  boxrule=0.5pt,
  before skip=2pt,
  after skip=2pt,
  top=1mm,bottom=1mm,left=2mm,right=2mm,
  title=\textbf{\color{primaryblue}#2},
  coltitle=primaryblue,
  attach boxed title to top left={yshift=-1.5mm,xshift=2.5mm},
  boxed title style={colback=lightgray,colframe=borderblue,arc=0.8mm,boxrule=0.5pt},
  #1
}


% --- OFFICIAL EXAM HEADER BOX ---
\newtcolorbox{headerbox}{
    colback=white,
    colframe=ForestGreen,
    boxrule=0.7pt,
    arc=1.5mm,
    left=2.5mm,
    right=2.5mm,
    top=1.5mm,
    bottom=1.5mm,
    before skip=1mm,
    after skip=1mm
}

\begin{document}


% --- OFFICIAL EXAM HEADER ---

\noindent
\begin{tabularx}{\textwidth}{@{}X >{\raggedleft\arraybackslash}p{6.2cm}@{}}
    \textbf{NATIONAL ECONOMICS UNIVERSITY} & \textbf{OFFICIAL MIDTERM EXAMINATION} \\
    \textbf{COLLEGE OF TECHNOLOGY} & {\small \shortstack[r]{Academic Year: 2026--2027 \\ Semester: Fall 2026}} \\
    \shortstack[l]{\textbf{FACULTY OF DATA SCIENCE \&} \\ \textbf{ARTIFICIAL INTELLIGENCE (FDA)}} & Course Code: \textbf{DSAI1003} \\
\end{tabularx}

\vspace{0.15em}

\begin{center}
    {\Large\bfseries\color{ForestGreen} FUNDAMENTAL PROGRAMMING CONCEPTS IN PYTHON}\\[0.10em]
    {\normalsize\textbf{Midterm Progress Exam: Part 01} \quad | \quad \textbf{Duration: 60 Minutes} (Excluding paper distribution)}
\end{center}

\vspace{0.05em}

\begin{headerbox}
\begin{tabularx}{\textwidth}{@{}X X@{}}
    \textbf{Candidate Full Name:} \dotfill & \textbf{Student ID (MSSV):} \dotfill \\[0.20em]
    \textbf{Class / Cohort:} \dotfill & \textbf{Exam Room / Seat No.:} \dotfill
\end{tabularx}

\vspace{0.20em}
\renewcommand{\arraystretch}{1.08}
\setlength{\tabcolsep}{5pt}
\centering
\begin{tabular}{|l|c|c|c|c||c|c|}
\hline
\textbf{Section} & \textbf{A} & \textbf{B} & \textbf{C} & \textbf{D} & \textbf{Total (100)} & \textbf{Examiner Signature} \\
\hline
\textbf{Max Score} & 20 & 25 & 25 & 30 & \textbf{100} & \\[0.25em]
\cline{1-6}
\textbf{Score} & & & & & & \\[0.38em]
\hline
\end{tabular}

\vspace{0.18em}
\raggedright\scriptsize
\textbf{Instructions:} Closed-book examination. All electronic devices and internet access are prohibited. Answer all questions directly in the designated areas. In Section D, use standard Python statements following the \textbf{Input $\rightarrow$ Process $\rightarrow$ Output} pattern; do \textbf{not} define custom functions (\texttt{def}) or use loops/branching. Do not detach pages.
\end{headerbox}

\vspace{0.6mm}

\noindent

\textbf{MULTIPLE-CHOICE ANSWER SHEET}
\textit{(Write A, B, C, or D in the corresponding box)}:

\begin{center}

\begin{tabular}{|c|c|c|c|c|c|c|c|c|c|c|}

  \hline

  \textbf{Question} & \textbf{1} & \textbf{2} & \textbf{3} & \textbf{4} &
  \textbf{5} & \textbf{6} & \textbf{7} & \textbf{8} & \textbf{9} & \textbf{10} \\

  \hline

  \textbf{Answer} & & & & & & & & & & \\[2mm]

  \hline

\end{tabular}

\end{center}

\vspace{-2.5mm}
\hrule
\vspace{0.8mm}

% ==============================================================================

% SECTION A

% ==============================================================================

\section*{Section A: Conceptual \& Architectural Fundamentals \hfill \normalsize\normalfont(20 points)}

\textit{Answer all 10 questions. Each correct answer carries 2 points. Suggested time: 12 minutes.}

\vspace{0.08cm}

\begin{enumerate}[label=\textbf{Q\arabic*.},leftmargin=2em,itemsep=2pt,topsep=1pt]

  \item Which computer hardware component is directly responsible for fetching, decoding, and executing program instructions?

  \begin{enumerate}[label=\Alph*.,noitemsep,topsep=0.5pt,leftmargin=2em]

    \item Primary Memory (RAM)

    \item Central Processing Unit (CPU)

    \item Secondary Storage (SSD / Hard Disk)

    \item Input/Output Bus

  \end{enumerate}

  \item What happens to the information stored in \textbf{Primary Memory (RAM)} when the computer is shut down?

  \begin{enumerate}[label=\Alph*.,noitemsep,topsep=0.5pt,leftmargin=2em]

    \item It is permanently stored in memory cells.

    \item It is lost because RAM is volatile memory.

    \item It is automatically compressed and transferred to cache.

    \item It is rewritten to CPU registers for immediate restart.

  \end{enumerate}

  \item Why can a single physical computer perform thousands of completely different tasks?

  \begin{enumerate}[label=\Alph*.,noitemsep,topsep=0.5pt,leftmargin=2em]

    \item The physical CPU hardware rewires itself automatically for each task.

    \item It can load and execute different software programs consisting of instructions and decisions.

    \item High-level languages bypass memory and directly command peripheral circuits.

    \item Secondary storage replaces CPU registers during complex computation.

  \end{enumerate}

  \item Which sequence accurately describes the execution pipeline of a standard Python program?

  \begin{enumerate}[label=\Alph*.,noitemsep,topsep=0.5pt,leftmargin=2em]

    \item Source Code (\texttt{.py}) $\rightarrow$ Machine Code $\rightarrow$ Compiler $\rightarrow$ Execution

    \item Source Code (\texttt{.py}) $\rightarrow$ Byte Code (\texttt{.pyc}) $\rightarrow$ Python Virtual Machine (PVM) $\rightarrow$ Running Program

    \item Byte Code (\texttt{.pyc}) $\rightarrow$ Source Code (\texttt{.py}) $\rightarrow$ Operating System $\rightarrow$ CPU

    \item Source Code (\texttt{.py}) $\rightarrow$ Linker $\rightarrow$ Binary File (\texttt{.exe}) $\rightarrow$ CPU

  \end{enumerate}

  \item Which of the following statements regarding high-level languages like Python is \textbf{FALSE}?

  \begin{enumerate}[label=\Alph*.,noitemsep,topsep=0.5pt,leftmargin=2em]

    \item Python code is portable across Windows, macOS, and Linux operating systems.

    \item The CPU directly executes Python source code text without needing an interpreter or virtual machine.

    \item High-level languages allow programmers to express algorithms using human-readable syntax.

    \item High-level languages free programmers from manually managing CPU registers and memory addresses.

  \end{enumerate}

  \newpage

  \item What are the \textbf{three essential characteristics} that every valid algorithm must satisfy?

  \begin{enumerate}[label=\Alph*.,noitemsep,topsep=0.5pt,leftmargin=2em]

    \item Fast, Portable, and Object-Oriented

    \item Unambiguous, Executable, and Terminating

    \item Compiled, Abstract, and Infinite

    \item Linear, Mathematical, and Formatted

  \end{enumerate}

  \item What is the main objective of creating \textbf{pseudocode} before writing Python code?

  \begin{enumerate}[label=\Alph*.,noitemsep,topsep=0.5pt,leftmargin=2em]

    \item To enable the Python interpreter to pre-compile variables into bytecode.

    \item To design and verify the algorithmic logic in plain language without getting bogged down by syntax details.

    \item To eliminate the need for primary memory during program execution.

    \item To generate automated unit tests for software applications.

  \end{enumerate}

  \item What data type does the built-in \texttt{input()} function \textbf{always} return in Python?

  \begin{enumerate}[label=\Alph*.,noitemsep,topsep=0.5pt,leftmargin=2em]

    \item \texttt{int}

    \item \texttt{float}

    \item \texttt{str}

    \item Depends automatically on whether the input text contains numeric digits.

  \end{enumerate}

  \item According to good programming practice, which sequence of problem-solving steps is most appropriate?

  \begin{enumerate}[label=\Alph*.,noitemsep,topsep=0.5pt,leftmargin=2em]

    \item Write Code $\rightarrow$ Run Program $\rightarrow$ Understand Problem $\rightarrow$ Pseudocode

    \item Understand Problem $\rightarrow$ Develop Algorithm/Pseudocode $\rightarrow$ Test by Hand $\rightarrow$ Translate to Python $\rightarrow$ Run \& Debug

    \item Design Algorithm $\rightarrow$ Understand Problem $\rightarrow$ Write Code $\rightarrow$ Test by Hand

    \item Write Code $\rightarrow$ Reverse-engineer Pseudocode $\rightarrow$ Debug

  \end{enumerate}

  \item Which of the following is the best example of a \textbf{Logic Error} (semantic error)?

  \begin{enumerate}[label=\Alph*.,noitemsep,topsep=0.5pt,leftmargin=2em]

    \item Omitting a closing parenthesis: \texttt{print("Hello"}.

    \item Attempting to divide a number by zero: \texttt{print(25 / 0)}.

    \item Calculating rectangle perimeter as \texttt{perimeter = length * width}\\
    instead of \texttt{2 * (length + width)}.

    \item Attempting to convert alphabetical letters to an integer: \texttt{int("twenty")}.

  \end{enumerate}

\end{enumerate}


% ==============================================================================

% SECTION B

% ==============================================================================

\section*{Section B: Code Tracing, Output Prediction \& Error Diagnosis \hfill \normalsize\normalfont(25 points)}

\textit{Answer all 5 questions. Each question carries 5 points. Suggested time: 15 minutes.}

\vspace{0.08cm}

\subsection*{Q11. Precise Output Prediction: \texttt{sep} and \texttt{end} (5 points)}

Carefully analyze the Python snippet below:

\begin{lstlisting}
print("2026", "09", "10", sep="-", end=" ")
print("EXAM", "PYTHON", sep="::")
print("DONE")
\end{lstlisting}

Write the \textbf{exact screen output} below. Clearly show any spaces and newlines:

\begin{answerbox}[height=2.0cm]{Output Box for Q11}

\vspace{0.03cm}

\end{answerbox}

\vspace{0.08cm}

\newpage

\subsection*{Q12. String Operations vs. Numeric Arithmetic (5 points)}

Given the following variable assignments and statements:

\begin{lstlisting}
a = "5"
b = "2"
c = 5
d = 2
print(a + b)
print(c + d)
print(a * 3)
print(c * 3)
\end{lstlisting}

Provide the exact printed output for each line:

\begin{itemize}[noitemsep]

  \item Output line 1 (\texttt{print(a + b)}): \underline{\hspace{4cm}}

  \item Output line 2 (\texttt{print(c + d)}): \underline{\hspace{4cm}}

  \item Output line 3 (\texttt{print(a * 3)}): \underline{\hspace{4cm}}

  \item Output line 4 (\texttt{print(c * 3)}): \underline{\hspace{4cm}}

\end{itemize}

\vspace{0.08cm}

\subsection*{Q13. Error Taxonomy and Classification (5 points)}

For each code segment, classify the error into \textbf{Compile-Time / Syntax Error}, \textbf{Exception (Run-time Error)}, or \textbf{Logic Error}. Provide a 1-sentence technical explanation.

\begin{center}

\renewcommand{\arraystretch}{1.6}

\begin{tabularx}{\textwidth}{|p{5.5cm}|p{4.2cm}|X|}

\hline

\textbf{Code Snippet} & \textbf{Error Classification} & \textbf{Technical Reason} \\

\hline

\textbf{A.} \texttt{print("Total is:" total)} & & \\[0.6cm]

\hline

\textbf{B.} \texttt{val = int("3.14159")} & & \\[0.6cm]

\hline

\textbf{C.} \texttt{avg = score1 + score2 + score3 / 3} \newline \textit{(Goal: calculate mean of 3 scores)} & & \\[0.8cm]

\hline

\end{tabularx}

\end{center}

\vspace{0.08cm}

\subsection*{Q14. Deconstructing Python Statements (5 points)}

Examine the following statement:

\begin{lstlisting}
# Display computed final score with bonus
print("Final Score:", base_score + 10)
\end{lstlisting}

Identify each corresponding element from the statement:

\begin{enumerate}[label=\alph*., noitemsep]

  \item The \textbf{comment}: \underline{\hspace{10cm}}

  \item The \textbf{built-in function name}: \underline{\hspace{8.2cm}}

  \item The \textbf{string literal argument}: \underline{\hspace{8cm}}

  \item The \textbf{arithmetic expression argument}: \underline{\hspace{7cm}}

  \item The \textbf{variable name}: \underline{\hspace{9.4cm}}

\end{enumerate}

\vspace{0.08cm}

\subsection*{Q15. Tracing Formatted f-Strings (5 points)}

Predict the exact screen output produced by the following snippet:

\begin{lstlisting}
item = "USB Drive"
quantity = 4
unit_price = 12.5
subtotal = quantity * unit_price
print(f"Item: {item} | Qty: {quantity}")
print(f"Subtotal: ${subtotal:.2f}")
print(f"Discount (10%): ${subtotal * 0.10:.2f}")
\end{lstlisting}

\begin{answerbox}[height=1.8cm]{Output Box for Q15}

\vspace{0.03cm}

\end{answerbox}


% ==============================================================================

% SECTION C

% ==============================================================================

\newpage

\section*{Section C: Algorithm Design \& Pseudocode \hfill \normalsize\normalfont(25 points)}

\textit{Answer all tasks. Suggested time: 15 minutes.}

\begin{tcolorbox}[colback=boxbg, colframe=secondaryteal, arc=1.5mm]

\textbf{Problem Specification: Electric Scooter Rental Billing System}\\[0.1cm]

A city electric scooter operator calculates trip fares based on the following rules:

\begin{itemize}[noitemsep]

  \item \textbf{Unlock Fee:} $\$1.50$ fixed rate per trip.

  \item \textbf{Ride Time Rate:} $\$0.25$ for each minute ridden.

  \item \textbf{Sales Tax:} An $8\%$ ($0.08$) municipal sales tax is added onto the subtotal fare (\text{Unlock Fee} + \text{Ride Cost}).

\end{itemize}

Your task is to design an algorithm that computes and outputs an itemized customer receipt.

\end{tcolorbox}

\vspace{0.08cm}

\noindent\textbf{Task 1: Identify Inputs and Outputs} (5 points)

\begin{itemize}[noitemsep]

  \item \textbf{Required Input(s):} \hrulefill \\[0.25cm]

  \item \textbf{Required Output(s):} \hrulefill

\end{itemize}

\vspace{0.08cm}

\noindent\textbf{Task 2: Write Structured Pseudocode} (10 points)\\

Write clear, sequential pseudocode adhering to the \textbf{Input $\rightarrow$ Process $\rightarrow$ Output} paradigm:

\begin{answerbox}[height=10.2cm]{Pseudocode Solution}

\vspace{0.03cm}

\end{answerbox}

\vspace{0.08cm}

\noindent\textbf{Task 3: Algorithm Verification} (5 points)\\

Demonstrate how your algorithm fulfills the three mandatory algorithm characteristics:

\begin{enumerate}[label=\alph*., noitemsep]

  \item \textbf{Unambiguous:} \hrulefill

  \item \textbf{Executable:} \hrulefill

  \item \textbf{Terminating:} \hrulefill

\end{enumerate}

\vspace{0.08cm}

\noindent\textbf{Task 4: Hand-Trace Calculation} (5 points)\\

Verify your algorithm manually for a trip duration of \textbf{20 minutes}:

\begin{itemize}[noitemsep]

  \item Unlock Fee: \underline{\hspace{3cm}} \hfill Riding Charge ($20 \times \$0.25$): \underline{\hspace{3cm}}

  \item Subtotal Before Tax: \underline{\hspace{3cm}} \hfill Tax Amount ($8\%$ of subtotal): \underline{\hspace{3cm}}

  \item \textbf{Final Total Fare to Pay:} \underline{\hspace{3.5cm}}

\end{itemize}


% ==============================================================================

% SECTION D

% ==============================================================================

\newpage

\section*{Section D: Applied Programming in Python \hfill \normalsize\normalfont(30 points)}

\textit{*Write complete, syntactically correct Python statements directly using the \textbf{Input $\rightarrow$ Process $\rightarrow$ Output} pattern. Do \textbf{not} define custom functions (\texttt{def}), \texttt{import}, or loops. Suggested time: 18 minutes.*}

\vspace{0.08cm}

\subsection*{Problem 1: Commuter Fuel Cost Calculator (15 points)}

Write a Python script that computes the fuel consumed and total monetary expense for a highway journey.

\begin{itemize}[noitemsep]

  \item Formulas:

  $$\text{fuel\_used} = \frac{\text{distance} \times \text{consumption\_rate}}{100}, \quad \text{trip\_cost} = \text{fuel\_used} \times \text{fuel\_price}$$

  \item Requirements:

  \begin{enumerate}[noitemsep]

    \item Prompt user for \texttt{distance} in km, \texttt{consumption\_rate} (L/100km), and \texttt{fuel\_price} per liter using \texttt{input()}. Ensure proper type conversion to \texttt{float}.

    \item Compute \texttt{fuel\_used} and \texttt{trip\_cost}.

    \item Display the summary using f-strings with values formatted to \textbf{2 decimal places (\texttt{:.2f})}.

  \end{enumerate}

\end{itemize}

\begin{answerbox}[height=13.2cm]{Python Solution for Problem 1}

\vspace{0.03cm}

\end{answerbox}

\vspace{0.12cm}

\newpage

\subsection*{Problem 2: Event Duration Breakdown (15 points)}

Write a Python script that converts a total number of elapsed seconds into \textbf{hours, minutes, and remaining seconds} using integer floor division (\texttt{//}) and modulus (\texttt{\%}).

\begin{itemize}[noitemsep]

  \item Requirements:

  \begin{enumerate}[noitemsep]

    \item Prompt user for \texttt{total\_seconds} and cast to \texttt{int}.

    \item Compute full hours ($1\text{ hr} = 3600\text{ s}$), remaining seconds, full minutes ($1\text{ min} = 60\text{ s}$), and leftover seconds.

    \item Display the result in the format: \texttt{"X hours, Y minutes, Z seconds"}.

  \end{enumerate}

\end{itemize}

\begin{answerbox}[height=16.0cm]{Python Solution for Problem 2}

\vspace{0.03cm}

\end{answerbox}

\begin{center}

  \vspace{0.15cm}

  \rule{0.5\textwidth}{0.4pt}\\[0.08cm]

  \textbf{\color{primaryblue}--- END OF EXAM PAPER ---}

\end{center}

\end{document}